% 图4-1：标注模型的准备与试标决策，按112 mm正文宽度设计。
\begin{tikzpicture}[
  x=1mm,
  y=1mm,
  font=\figurefont\footnotesize,
  text=BookInk,
  prep input/.style={book node,draw=FigureInput,fill=FigureInput!7,line width=0.8pt,text width=25mm,minimum height=11mm,inner sep=1.5mm},
  prep process/.style={book node,draw=FigureModel,fill=FigureModel!7,line width=0.8pt,text width=28mm,minimum height=11mm,inner sep=1.5mm},
  prep decision/.style={book node,draw=BookAmber,fill=BookAmber!7,line width=0.9pt,text width=24mm,minimum height=11mm,inner sep=1.5mm},
  prep result/.style={book node,draw=BookMuted,fill=BookPaper,line width=0.7pt,text width=25mm,minimum height=11mm,inner sep=1.5mm},
  prep flow/.style={-{Stealth[length=2mm]},draw=BookMuted,line width=0.8pt,rounded corners=1mm},
  prep feedback/.style={-{Stealth[length=2mm]},draw=BookAmber,line width=0.8pt,dashed,rounded corners=1mm},
  edge text/.style={fill=white,inner sep=0.8pt,font=\figurefont\scriptsize,text=BookMuted}
]
  \node[prep input] (requirements) at (16,0) {任务要求\\类别、边界、弃权};
  \node[prep input] (models) at (55,0) {可用模型\\专用或通用};
  \node[prep input] (labels) at (94,0) {初始标签\\数据与覆盖};

  \node[prep process] (select) at (55,-19) {选择候选方案\\并记录资源约束};
  \node[prep decision] (need) at (55,-37) {是否需要\\参数训练？};
  \node[prep process] (adapt) at (87,-51) {训练或适配\\形成模型版本};
  \node[prep process] (pilot) at (55,-66) {试标确认\\能力范围与失败类型};

  \node[prep result,draw=BookAmber] (source) at (17,-86) {补充其他来源\\或重定任务};
  \node[prep result,draw=FigureOutput] (accept) at (55,-86) {采用当前方案\\进入批量获取};
  \node[prep result,draw=FigureModel] (continue) at (93,-86) {继续适配\\再做试标};

  \draw[prep flow] (requirements.south) -- (16,-10) -| (select.west);
  \draw[prep flow] (models.south) -- (select.north);
  \draw[prep flow] (labels.south) -- (94,-10) -| (select.east);
  \draw[prep flow] (select.south) -- (need.north);
  \draw[prep flow] (need.south) -- node[edge text,right] {否} (pilot.north);
  \draw[prep flow] (need.east) -- node[edge text,above] {是} (72,-37) |- (adapt.north);
  \draw[prep flow] (adapt.south) |- (pilot.east);
  \draw[prep flow] (pilot.south) -- (accept.north);
  \draw[prep flow] (pilot.west) -- (36,-66) |- (source.north);
  \draw[prep flow] (pilot.east) -- (75,-66) |- (continue.north);

  \draw[prep feedback] (source.west) -- (1,-86) |- (select.west);
  \draw[prep feedback] (continue.east) -- (110,-86) |- (adapt.east);
  \node[edge text,anchor=west] at (3,-29) {资源或要求变化};
  \node[edge text,anchor=east] at (108,-42) {新训练版本};
\end{tikzpicture}
